\section{Related Work}
\label{sec:related_work}

Our proposed framework, ELATI, intersects several prominent domains in modern machine learning research: parameter-efficient multi-tenant serving, weight-space model merging, dynamic routing, and the analysis of representation dynamics in deep networks. Below, we contextualize our work relative to these areas.

\subsection{Static Model Merging and Weight Interpolation}
Weight-space model merging has emerged as a compelling paradigm for combining multiple specialized networks without retraining.
Early efforts demonstrated that averaging the weights of models trained from the same initialization (e.g., Model Soups \cite{wortsman22soups}) can improve generalizability.
To address multi-task adaptation, Task Arithmetic \cite{ilharco22arithmetic} introduced the addition and subtraction of task-specific weight vectors (fine-tuned weights minus pre-trained weights) directly on the base model.
Subsequent works identified that static merging suffers from severe ``representation conflicts'' where parameter updates for different tasks interfere and cancel each other out, collapsing task-specific performance.
To resolve this, TIES-Merging \cite{yadav23ties} filters out redundant parameter updates, resolves sign agreements, and merges only dominant directions.
DARE \cite{yu23dare} prunes insignificant parameter updates and rescales them, achieving high performance with extremely sparse updates.
Other approaches explore advanced weight scaling, such as Fisher weighted averaging \cite{matena22fisher}, regression-based alignment (RegMean \cite{jin22regmean}), permutation symmetry alignment (Git Re-Basin \cite{ainsworth22gitrebasin}), and alignment of disparate architectures (ZipIt! \cite{stoica23zipit}).
Despite these innovations, static model merging remains fundamentally limited: a single, static set of merged weights represents a uniform trade-off that cannot adapt to the sample-specific characteristics of a dynamic, multi-task deployment stream on-the-fly.

\subsection{Mixture of Experts and Parameter-Efficient serving}
To scale model capacity dynamically, Mixture of Experts (MoE) \cite{shazeer17moe, wang24moesurvey} routes samples to specialized subnetworks (experts) using a parameterized gating network.
While full MoE architectures like Switch Transformers \cite{fedus21switch} and GShard \cite{lepikhin20gshard} offer massive capacity, they require immense VRAM and specialized hardware clusters, making them inaccessible for resource-constrained edge systems.
To democratize multi-task serving, researchers have combined parameter-efficient fine-tuning (PEFT), such as LoRA \cite{hu21lora, mangrulkar22peft}, with MoE structures.
MoE-LoRA \cite{dou23moelora} and MoE-LLaVA \cite{liu24moellava} assign low-rank adapters as experts under a routing gate.
Similarly, LoraHub \cite{huang23lorahub} proposes zero-shot composition of LoRA modules based on user instructions.
While PEFT-based MoE models significantly reduce VRAM requirements, they still require executing parallel forward paths through different adapter paths or storing multiple active expert copies in memory during batch execution, which incurs substantial routing latency and memory footprint on standard hardware.

\subsection{Dynamic Routing and Model Merging}
Rather than routing activations through physical subnetworks, dynamic model merging interpolates task-specific adapter weights into the base model on-the-fly.
This preserves a single-model VRAM footprint while dynamically customizing the model for the incoming batch.
Parameter-Free Subspace Routing (PFSR) combined with Micro-Batch Homogenization (MBH) represents the state-of-the-art in this direction \cite{chronopoulou23representation, zhang24survey}.
PFSR calculates routing coefficients by projecting intermediate representations against the expert classification heads, while MBH groups heterogeneous samples into homogeneous micro-batches to perform dynamic parameter interpolation.
However, extracting representations from the penultimate layer of the backbone introduces a massive \emph{two-pass latency penalty}.
The system must run a complete, throw-away forward pass of the deep base model backbone just to obtain routing coefficients, and then run a second pass through the dynamically merged network.
This latency penalty makes penultimate dynamic merging highly impractical for high-throughput, low-latency deployments.
ELATI directly addresses this limitation, achieving a true one-pass dynamic merging pipeline.

\subsection{Layer-wise Representation Dynamics and Early Probing}
Our work is also grounded in the study of representation learning and BERTology \cite{rogers20bertology}.
Linguistic and semantic probing studies (e.g., \cite{tenney19pipeline, liu19linguistic, belinkov22probing}) demonstrate that deep networks construct representations hierarchically.
While late layers contain task-specific semantic outputs, early layers capture syntactic, structural, and domain-level characteristics.
Linear classifier probes \cite{alain16probing, conneau18probe} have shown that intermediate layer representations can successfully predict downstream task labels long before the final layer.
This hierarchical representation is exploited by early-exit networks (such as BranchyNet \cite{teerapittayanon16branchynet}, Shallow-Deep Networks \cite{kaya19sdn}, DeeBERT \cite{xin20deebert}, and CALM \cite{schuster22calm}) to bypass deep layers for easy samples.
Unlike early-exit networks that terminate computation to save FLOPs, ELATI leverages early-layer representations to identify the target task, dynamically merging and executing the \emph{remaining} downstream layers on-the-fly.
To our knowledge, ELATI is the first framework to utilize unsupervised early-layer representation centroids to drive training-free, parameter-free dynamic weight-space merging.
